\begin{abstract}
Large language models often change a correct answer once a user pushes back, even without new evidence. We test whether a model's uncertainty before that pressure begins predicts this failure. Using six models from the GPT, Claude, and Gemini families on MMLU-Pro, GPQA-Diamond, and Humanity's Last Exam, we estimate each model's pre-pressure belief by repeated sampling, then apply a fixed sequence of six escalating pressure turns asserting a specific wrong answer. We find that questions where the model starts out uncertain flip to the wrong answer at close to twice the rate of questions where it starts out certain. Pressure level and pre-pressure uncertainty also interact, significantly for only one model in isolation, though pooling across models with a meta-analysis recovers a small but significant effect. Across six targets, models split into three patterns under sustained pressure: some resist, some capitulate further with each turn, and some stop reaching a consistent answer across attempts. The pattern depends on both the model and benchmark difficulty, not on a fixed model property. In a chain-of-thought variant, belief drifts to the wrong answer well before confidence drops, so confidence alone is a poor signal for catching this failure. Susceptibility to sycophancy is dependent on how uncertain the model already was.
\end{abstract}
